% ID: SEL-D0019-OB-E04
% ESTADO: cotejada. Segunda pasada visual completa y borrador compilado revisados.
% Fuente: Selectividad/matematicas_aplicadas/reserva a/Examen-Matemáticas Aplicadas a las Ciencias Sociales.pdf
% SHA-256: 6fed0384ff6fcbaaa7098ab3f35ea5f881d1280502f050ada3a7f13af528392e
% Página: 2; opción B; ejercicio original 4.

Se desea estimar el porcentaje de jóvenes que utilizan una determinada red social. Para ello se escoge una muestra aleatoria simple de 500 jóvenes y de ellos 410 afirman utilizarla.
\begin{enumerate}[label=\alph*)]
\item \textbf{(1.5 puntos)} Calcule el intervalo de confianza para la proporción de jóvenes que usa esa red social con un nivel de confianza del 95\%.
\item \textbf{(1 punto)} Manteniendo la proporción muestral, determine el tamaño mínimo de la muestra necesario para que, con un nivel de confianza del 97\%, el error máximo que se cometa al estimar la proporción de esa población sea inferior a $0.04$.
\end{enumerate}
